% !TEX root = ../main.tex
%% Devoir surveillé n° 3 — chapitres 8 à 10 (valeurs vérifiées : verifications/devoirs.py)

\devoirDS{3}{Chapitres 8 à 10}{{L'arôme de banane : synthèse, rendement et spectroscopie IR}/8,{Le réchaud de camping : combustion et énergie}/7,{Le bioéthanol : molécule et bilan carbone}/5}

\begin{exoBac}{8}{L'arôme de banane}
L'éthanoate de 3-méthylbutyle est l'ester responsable de l'odeur de banane. On le synthétise à partir d'acide éthanoïque et de
3-méthylbutan-1-ol, en présence de quelques gouttes d'acide sulfurique, selon une transformation lente et non totale :
\[\ce{CH3-COOH + (CH3)2CH-CH2-CH2-OH -> CH3-COO-CH2-CH2-CH(CH3)2 + H2O}\]
\begin{docu}{1}{Le protocole}
Dans un ballon, on introduit \qty{21,7}{mL} de 3-méthylbutan-1-ol, \qty{23,0}{mL} d'acide éthanoïque, quelques gouttes d'acide sulfurique et
de la pierre ponce, puis on chauffe à reflux pendant une heure. Après lavage, séchage et distillation, on obtient \qty{20,8}{g} d'ester.
\end{docu}
\begin{donnees}
Masses molaires : acide éthanoïque \qty{60,0}{g.mol^{-1}} ; 3-méthylbutan-1-ol \qty{88,0}{g.mol^{-1}} ; ester \qty{130,0}{g.mol^{-1}}. Densités :
acide éthanoïque 1,05 ; 3-méthylbutan-1-ol 0,81 ; ester 0,87. Bandes infrarouges : O--H d'un alcool 3200 -- \qty{3400}{cm^{-1}}, large ; O--H d'un acide 2500 -- \qty{3200}{cm^{-1}}, très large ;
C--H 2800 -- \qty{3100}{cm^{-1}} ; C$=$O d'un ester 1735 -- \qty{1750}{cm^{-1}} ; C$=$O d'un acide 1700 -- \qty{1725}{cm^{-1}}.
\end{donnees}
\begin{docu}{2}{Deux spectres infrarouges}
\schema{ds3-ir}
\end{docu}
\begin{questions}
\item Recopier les formules semi-développées des trois espèces organiques de l'équation, entourer leurs groupes caractéristiques
et nommer les familles correspondantes.
\item Faire un schéma légendé du montage de chauffage à reflux et préciser son intérêt.
\item Calculer les quantités de matière initiales des réactifs et identifier le réactif limitant.
\item Calculer la masse maximale d'ester que l'on pourrait obtenir, puis le rendement de la synthèse.
\item Lequel des deux spectres est celui de l'ester obtenu ? Justifier.
\item Pourquoi introduit-on l'acide éthanoïque en excès ?
\end{questions}
\end{exoBac}

\begin{exoBac}{7}{Le réchaud de camping}
Un réchaud de camping brûle le butane $\ce{C4H10}$ contenu dans une cartouche de \qty{190}{g}. La combustion est supposée complète : elle
ne produit que du dioxyde de carbone et de l'eau.
\begin{donnees}
Énergies de liaison (\unit{kJ.mol^{-1}}) : C--C 345 ; C--H 415 ; O$=$O 498 ; C$=$O (dans $\ce{CO2}$) 804 ; O--H 463. Masses molaires :
$M(\ce{C4H10})=\qty{58,0}{g.mol^{-1}}$ ; $M(\ce{CO2})=\qty{44,0}{g.mol^{-1}}$. Pouvoir calorifique du butane (tables) : \qty{45,8}{MJ.kg^{-1}}. Il faut
\qty{4,18}{kJ} pour élever de \qty{1}{\degreeCelsius} la température de \qty{1}{kg} d'eau ; \qty{1,0}{L} d'eau a une masse de \qty{1,0}{kg}. Seuls
\qty{40}{\percent} de l'énergie libérée par la combustion chauffent l'eau de la casserole.
\end{donnees}
\begin{questions}
\item Écrire l'équation de la réaction de combustion complète du butane.
\item À l'aide des énergies de liaison, estimer l'énergie molaire de cette combustion, en phase gazeuse, pour une mole de butane.
\item En déduire une estimation du pouvoir calorifique du butane et la comparer à la valeur des tables.
\item Quelle énergie la combustion du contenu d'une cartouche libère-t-elle ?
\item \initiative Combien de litres d'eau peut-on porter de \qty{20}{\degreeCelsius} à \qty{100}{\degreeCelsius} avec une cartouche ?
\item Quelle masse de dioxyde de carbone la combustion d'une cartouche rejette-t-elle ?
\end{questions}
\end{exoBac}

\begin{exoBac}{5}{Le bioéthanol}
Le bioéthanol, produit par fermentation de betteraves ou de céréales, est mélangé à l'essence de certains carburants. C'est de
l'éthanol, de formule brute $\ce{C2H6O}$.
\begin{donnees}
Masse molaire de l'éthanol : \qty{46,0}{g.mol^{-1}} ; son pouvoir calorifique : \qty{26,8}{MJ.kg^{-1}}. Masses molaires : $M(\ce{C4H10})=\qty{58,0}{g.mol^{-1}}$ ;
$M(\ce{CO2})=\qty{44,0}{g.mol^{-1}}$. Pouvoir calorifique du butane : \qty{45,8}{MJ.kg^{-1}}.
\end{donnees}
\begin{questions}
\item Écrire la formule semi-développée de l'éthanol. Nommer son groupe caractéristique et sa famille.
\item Le méthoxyméthane, $\ce{CH3-O-CH3}$, a la même formule brute. Comment la spectroscopie infrarouge permet-elle de distinguer ces
deux espèces ?
\item Écrire l'équation de la combustion complète de l'éthanol.
\item Calculer la masse de dioxyde de carbone rejetée pour produire \qty{1}{MJ} d'énergie, en brûlant de l'éthanol, puis du butane. Comparer.
\item Pourquoi considère-t-on pourtant que le bioéthanol contribue moins à l'effet de serre ?
\end{questions}
\end{exoBac}
